\documentclass{report}
\usepackage[T1]{fontenc} 
\usepackage[utf8]{inputenc} 
\usepackage{lmodern}
\usepackage [frenchb]{babel}
% Pour pouvoir utiliser 
\usepackage{ucs}

\usepackage{textcomp}
\usepackage{graphicx}
\usepackage{keystroke}
\usepackage{amssymb}
\usepackage{amsmath}
%\renewcommand{\thesection}{\arabic{section}} % numérotation des sections	
\usepackage[cc]{titlepic} %rajouter le logo dans la page de garde
\usepackage{url} % Pour avoir de belles url
\usepackage[a4paper]{geometry}
\usepackage[linktocpage]{hyperref}
% Pour pouvoir faire plusieurs colonnes
\usepackage {multicol}
%Pour faire plusieurs lignes
\usepackage{multirow}
% Pour mettre du code source
\usepackage {listings}
% Pour pouvoir passer en paysage
\usepackage{lscape}	

% Pour pouvoir faire plusieurs colonnes
\usepackage {multicol}

% POur crééer un index
\usepackage{makeidx}

%insertion de pdf
\usepackage{pdfpages}


%Pour taper des algorithme en pseudo code
%\usepackage{algorithm,algorithmic}
%Utilisation d'acronyme

\usepackage{fancyhdr}
\pagestyle{fancy}
\usepackage{lastpage} %numérotation type k/n
\renewcommand{\footrulewidth}{0.5pt} % trait horizontal
%\renewcommand{\headrulewidth}{1pt} % suppresion du trait horizontal dans l'entête
\lfoot{Senior Capstone Project}
\cfoot{\hyperref[tableofcontents]{\thepage} sur \pageref*{LastPage} } % numérotation des pages
\rfoot{Master ALMA}
\fancyhead[L]{\hspace{-1cm}\includegraphics[scale=0.105]{img/logouniv.png}} %logo entête !


\fancypagestyle{plain}{% 1ères pages des chapitres:
  % \fancyhead{} % supprime l’entete...
  % \renewcommand{\headrulewidth}{0pt} % ...et le filet
}



\hypersetup{
  backref=true,
  %permet d'ajouter des liens dans...
  pagebackref=true,%...les bibliographies
  hyperindex=true, %ajoute des liens dans les index.
  colorlinks=true, %colorise les liens
  breaklinks=true, %permet le retour à la ligne dans les liens trop longs
  urlcolor= blue, %couleur des hyperliens
  citecolor=	cyan,
  bookmarks=true, %créé des signets pour Acrobat
  bookmarksopen=true,
  %si les signets Acrobat sont créés,
  %les afficher complètement.
  pdftitle={Rapport Senior Capstone Project}, %informations apparaissant dans
  pdfauthor={Jonathan BIZET, Jeremy CHALINE, Thomas DEUGUEULE, Algassimou DIALO, Pierre GUILLET, Romain RINCE},
  %dans les informations du document
  %sous Acrobat.
}

\makeindex
%%%% debut macro pour enlever le nom chapitre %%%%
\makeatletter
\def\@makechapterhead#1{%
  \vspace*{50\p@}%
          {\parindent \z@ \raggedright \normalfont
            \interlinepenalty\@M
            \ifnum \c@secnumdepth >\m@ne
            \Huge\bfseries \thechapter\quad
            \fi
            \Huge \bfseries #1\par\nobreak
            \vskip 40\p@
}}

\def\@makeschapterhead#1{%
  \vspace*{50\p@}%
          {\parindent \z@ \raggedright
            \normalfont
            \interlinepenalty\@M
            \Huge \bfseries  #1\par\nobreak
            \vskip 40\p@
}}
\makeatother
%%%% fin macro %%%%
%Couverture 
\widowpenalty=10000
\clubpenalty=10000
\hyphenation{permet-tant}


%\vspace{20cm}
\newcommand{\HRule}{\rule{\linewidth}{0.5mm}}

\begin{document}
%%%%%%%%%%%%%%%%%%%%%%%%%%%%%%%%%%%%%%%%%
%		Page de garde		%
%%%%%%%%%%%%%%%%%%%%%%%%%%%%%%%%%%%%%%%%%

\begin{titlepage}
\pagestyle{empty}
\begin{center}
% Upper part of the page
% Title
\textsc{\LARGE Master 2\\[0.15cm]}\emph{\LARGE Architecture Logicielle}\\
\hfill\\[1.5cm]
\HRule \\[0.4cm]
\textsc{\huge Senior Capstone Project}\\[0.5cm]
{\Large \emph{Injection de tests dans des composants logiciels par transformation de modèles}}\\[0.4cm]

\HRule \\[.5cm]

% Author and supervisor
\begin{minipage}{0.4\textwidth}
\begin{flushleft} \large
\emph{Auteurs :}\\
Jonathan \textsc{Bizet}\\Jeremy \textsc{Chaline}\\Thomas \textsc{Degueule}\\Algassimou \textsc{Diallo}\\Pierre \textsc{Guillet}\\Romain \textsc{Rincé}
\end{flushleft}
\end{minipage}
\begin{minipage}{0.4\textwidth}
\begin{flushright} \large
\emph{Encadrants :} \\
P.\textsc{André}\\G.\textsc{Ardourel}\\J.M.\textsc{Mottu}
\end{flushright}
\end{minipage}\\[4cm]

\includegraphics[width=0.35\textwidth]{img/logouniv}

\vfill
\hfill



{\large \today}\\[0.5cm]
\emph{  Université de Nantes\\
2 rue de la Houssinière, BP92208, F-44322 \textsc{Nantes} cedex 03, \textsc{France}}\\
% Bottom of the page

\end{center}

\end{titlepage}
\clearpage
%%%%%%%%%%%%%%%%%%%%%%%%%%%%%%%%%
%	Fin page de garde	%
%%%%%%%%%%%%%%%%%%%%%%%%%%%%%%%%%

\renewcommand{\labelitemi}{$\bullet$} 

\clearpage


\tableofcontents\label{tableofcontents}
\clearpage

\chapter{Introduction}\label{chap:Intro}

\section{Problématiques}
Les problèmes d'inter-opérabilité et de réutilisabilité font partie des problématiques majeures des systèmes informatiques. La programmation objet ne parvient pas à répondre à ces problématiques à cause notamment de ses propriétés de granularité et de fort couplage. Les problèmes de cohésions (pas d'explications des besoins requis) sont un frein à l'utilisation de la programmation objet à grande échelle. 

Parmi les solutions à ces problématiques les architectures à composants sont des alternatives efficaces. Un composant, entité principale, correspond par définition à une ou plusieurs fonctionnalités. Il est donc possible de le concevoir et de le pré-tester de manière indépendante. La coopération entres composant est alors plus aisée. De plus un composant possède un couplage lâche. Cette propriété garantissant l'inter-opératibilté est permise par l'utilisation d'interfaces entre composants. 

\section{Présentation du sujet}
Le module \textsc{Architectures and Components} nous propose une étude de la programmation par composants. L'objectif de ce Projet/TP  est de réaliser un Home Architecture Description Language, et d'implémenter une architecture classique \emph{Clients-Serveur}. La démarche proposée pour atteindre cet objectif est composée de trois étape : 


\begin{enumerate}
\item 
Réalisation du M2 :\hfill \\
Il s'agit d'un méta-modèle de l'architecture en composants.  
\item
Réalisation du M1 :\hfill \\
À partir du méta-modèle précédement réalisé, nous pouvons proposer une instance de ce dernier et obtenir un modèle. Nous choisirons ici de modéliser une application classique \emph{Client-Serveurs}.  
\item
Réalisation du M0 :\hfill \\
Dernière étape, elle consiste à l'implémentation du M1 et de coder concrétement qui y sont modélisées.
\end{enumerate}

Dans ce rapport nous détaillerons chaque étape la modélisation proposée puis l'implémentation.

\input{Contexte}
\chapter{Analyse du problème}

\input{Proposition}
\input{Resultat}
\input{Gestion}
\chapter{Conclusion}\label{chap:COnc}
À ce stade, la modélisation de notre méta-modèle (\cf{} chapitre \ref{chap:M2}) nous a permis de générer un modèle pour notre application \textit{Clients-Serveur} (\cf{} chapitre \ref{chap:M1}). Par la suite il sera nécessaire de passer à l'étape 3 présentée au chapitre \ref{chap:Intro}. C'est à dire  d'implémenter le méta-modèle M2 en JAVA.

La création de l'application \textit{Clients-Serveur} repose sur l'héritage des classes du méta-modèle. Enfin l'instanciation concrète de ses classes achèvera l'implémentation de notre application.

% LocalWords:  méta-modèle Clients-Serveur l'instanciation
% LocalWords:  l'implémentation




\listoffigures
% Pour finir l'interligne de 1,5
%\bibliographystyle{alpha}
%\bibliography{biblio.bib}


\end{document}
